\chapter{Limitações}
\label{ch:limitacoes}

\section{Condições de uso}
\label{sec:funcao-limitacoes}

Os resultados do \CabotageLens{} são estimativas condicionadas por rota, carga, portos, fontes de distância, cobertura ANTAQ/MRV, regras de alocação e fronteiras ambiental e econômica. As limitações apresentadas neste capítulo definem o alcance das conclusões e indicam quais evidências adicionais seriam necessárias para uma decisão comercial ou operacional.

\begin{table}[htbp]
\centering
\small
\caption{Fronteiras que delimitam o uso dos resultados.}
\label{tab:condicoes-uso}
\begin{tabularx}{\textwidth}{L{0.23\textwidth}Y}
\toprule
Fronteira & Condição de interpretação \\
\midrule
Dados marítimos & Viagens ANTAQ representam operações registradas; intensidades EU MRV são proxies técnicos com cobertura variável por IMO.\\
Estatística & A mediana ponderada representa o centro da distribuição segundo o trabalho de transporte; não elimina viés de amostra ou de fallback.\\
Rota e porto & Corredor observado, porto selecionado, distância marítima e casos \texttt{same-port} condicionam a força do cenário.\\
Ambiental & Emissões operacionais \ttw{} de \cooe{}, sem inferir \wtw{}, LCA, CO$_2$ isolado ou inventário completo de poluentes.\\
Econômica & Custos modelados, não fretes comerciais, tarifas, contratos ou recomendação de compra.\\
Operacional & Rotas modeladas não comprovam serviço, frequência, slot, terminal, agenda ou aceitação de carga.\\
Validação & Sensibilidades e benchmark sustentam interpretações específicas; concordância direcional não equivale a calibração de magnitude.\\
\bottomrule
\end{tabularx}
\end{table}

\section{Dados ANTAQ, reconstrução da viagem e carga a bordo}
\label{sec:limitacao-dados-maritimos}

A sequência dos portos é reconstruída a partir das atracações registradas pela ANTAQ. A qualidade do corredor depende da completude dos identificadores de viagem, dos horários, da normalização dos portos e da associação entre atracações e movimentações de carga. Ausências, duplicidades ou inconsistências nesses campos podem alterar a ordem das escalas ou impedir a utilização de uma observação.

A carga a bordo é inferida a partir dos embarques e desembarques vinculados às escalas. A reconstrução adota o menor carregamento inicial não negativo capaz de compatibilizar a sequência de movimentações. Essa regra garante consistência aritmética, mas não observa diretamente o manifesto completo do navio, cargas fora do recorte, reposicionamento de contêineres vazios ou correções operacionais não refletidas nas tabelas.

Uma viagem entra na amostra de um par OD quando a origem precede o destino. Essa regra identifica trajetórias compatíveis com a ligação analisada, mas não comprova que a mesma carga permaneceu a bordo durante todo o percurso nem que o serviço continua comercialmente disponível. Os registros representam operações observadas no período das bases utilizadas.

\section{Intensidade MRV, fallbacks e outliers}
\label{sec:limitacao-intensidade-mrv}

O EU MRV fornece indicadores de navios que reportam no contexto regulatório europeu. A correspondência por IMO melhora a especificidade da estimativa, mas o indicador continua sendo um proxy de desempenho: condições brasileiras de velocidade, calado, mar, corrente, manutenção, combustível e operação podem diferir das condições consolidadas no MRV.

Os valores exatos de IMO são preservados sem aparagem. Essa decisão mantém a observação individual auditável, inclusive quando sua intensidade é elevada, mas não garante que todo valor extremo esteja livre de erro de origem ou seja representativo da viagem ANTAQ. A mediana ponderada no nível OD reduz a influência desses extremos sobre o indicador final; os valores individuais e a média aritmética diagnóstica permanecem disponíveis para inspeção.

Os fallbacks por classe ou tipo são estatísticas de grupo. A média aparada de 1\% ou a mediana protege contra caudas extremas, mas não substitui uma medição do navio ausente. Diferenças de porte, idade, motor, perfil operacional e qualidade da classificação permanecem agregadas dentro do grupo. Em Santos--Manaus, 49 das 89 observações usam fallback e representam 59,54\% do trabalho de transporte; por isso, a intensidade de 9,322050 g/(t$\cdot$nm) deve ser lida junto com essa composição de fonte.

A ponderação pelo trabalho de transporte atribui maior influência a viagens com maior produto carga--distância. Essa escolha é coerente com a unidade da intensidade, mas faz com que erros na reconstrução da carga ou na distância afetem simultaneamente o peso estatístico e o consumo calculado. A mediana ponderada limita a influência de intensidades extremas; ela não corrige erros sistemáticos compartilhados por grande parcela do trabalho de transporte.

\section{Fronteira de rota: portos, corredores e distâncias}
\label{sec:limitacao-rota}

O corredor selecionado representa a geometria do cenário, não uma previsão de itinerário. A prioridade dada à ligação direta observada e, na sua ausência, ao corredor de menor distância é uma regra determinística. Ela não considera frequência, tempo de espera, transbordo, capacidade, confiabilidade, preço comercial ou preferência do embarcador.

Uma distância proveniente da matriz marítima ou de corredor observado oferece suporte mais forte do que \texttt{haversine\_fallback}. A distância de grande círculo é útil para manter um cenário identificável quando falta uma rota navegável, mas não representa canais, desvios costeiros, restrições hidrográficas ou separação de tráfego.

Portos alternativos não são equivalentes por proximidade regional. Pecém pode compor um cenário para a região de Fortaleza, mas não valida o Porto de Fortaleza; Suape pode compor um cenário para a região de Recife, mas não valida o Porto do Recife. A troca modifica acessos rodoviários, distância marítima, terminal e condições de serviço. Casos \texttt{same-port}, como Santos/Santos, não contêm perna marítima substantiva entre portos distintos e permanecem fora da comparação modal principal.

\section{Fronteira ambiental}
\label{sec:limitacao-ambiental}

O \CabotageLens{} reporta emissões operacionais \ttw{} de \cooe{} associadas aos componentes representados. Para as pernas rodoviárias, a fronteira corresponde à combustão no veículo; para a perna marítima, à combustão a bordo e às parcelas operacionais explicitamente incluídas.

Produção e distribuição de combustíveis, fabricação e manutenção de veículos e navios, construção de infraestrutura, fim de vida, combustíveis alternativos e emissões locais completas pertencem a fronteiras WTT, \wtw{} ou LCA. Fontes que reportam CO$_2$ isolado ou \cooe{} sob outro conjunto de gases exigem reconciliação explícita antes de comparação direta \cite{icct2022,decarb2024,maritimelca2024}.

Operações portuárias são componentes modelados, não medições completas de terminais específicos. Quando a intensidade MRV por trabalho de transporte é aplicada, o \emph{hoteling} não é somado separadamente, sob a hipótese de que o indicador anual abrange o consumo reportado do navio para essa fronteira. Essa escolha reduz o risco de dupla contagem, mas não estima dispersão atmosférica, exposição humana, produtividade do berço ou qualidade do ar local \cite{berth2009,shipops2022,berthairquality2010}.

\section{Fronteira econômica}
\label{sec:limitacao-economica}

Os valores em BRL constituem um proxy de custo operacional modelado. A alternativa rodoviária representa um veículo dedicado à carga do cenário, enquanto a perna marítima distribui custo e consumo segundo a alocação definida. Essa assimetria operacional faz parte da unidade funcional adotada, mas pode produzir diferenças maiores do que as observadas em cotações comerciais.

Frete negociado, margem, seguro, demurrage, detention, inventário em trânsito, variação de prazo, risco, taxas locais, tarifas de terminal e disponibilidade de slot não estão integralmente representados. Um resultado favorável em custo indica uma hipótese para investigação, não uma conclusão de viabilidade comercial \cite{competitiveness2024,modalshiftreview2020}.

\section{Fronteira operacional}
\label{sec:limitacao-operacional}

O modelo também não resolve uma super-rede multimodal com operadores, transbordos, frequências, capacidades, tempos de espera e serviços concorrentes. Porto e corredor funcionam como elementos do cenário determinístico. A execução real requer confirmação de linha, janela, terminal, documentação, capacidade, aceitação da carga e coordenação entre operadores.

\section{Validação}
\label{sec:limitacao-validacao}

O benchmark externo oferece concordância direcional em 21 pares, mas não reprodução de magnitude. A classificação \codecell{same\_direction\_large\_gap} limita seu uso a uma verificação qualitativa da direção modal. Distâncias, carga, alocação, portos e fronteira ambiental precisam ser reconciliados antes de qualquer comparação quantitativa.

\section{Fontes e generalização}
\label{sec:limitacao-fontes-generalizacao}

A literatura contextualiza cabotagem, mudança modal, fronteiras ambientais e barreiras comerciais. Valores de outros estudos não são convertidos automaticamente em coeficientes do \CabotageLens{}. Qualquer incorporação exige compatibilidade de unidade funcional, carga, rota, fronteira, fonte e implementação.

Por fim, os resultados são específicos aos cenários analisados. Sua generalização para outros corredores, cargas, navios, portos ou períodos depende de cobertura de dados equivalente e de nova verificação das premissas. A ferramenta sustenta triagem acadêmica e comparação auditável; não sustenta, por si só, superioridade universal da cabotagem.
